\begin{frame}{단일 시드 vs 20회 평균}
  \begin{columns}[T]
    \begin{column}{0.56\textwidth}
      \centering
      \includegraphics[width=\linewidth]{ch07_2_vconverge.png}
      {\footnotesize 단일 실행(시드 0, $\alpha=0.2$, 10 에피소드)의
      학습 $V(s)$와 참값 $V^*=s/6$. $n=1$(회색)은 거의 0에서
      움직이지 못해 큰 편향, $n=5$(파랑)은 중간, $n=50$(빨강)은
      이번 실행에선 참값에 가장 가까움.}
    \end{column}
    \begin{column}{0.42\textwidth}
      \begin{block}{함정의 본질}
        \kq{하이퍼파라미터를 고르면서 단일 실행(단일 시드, 단일
        검증집합)에 의존하면, ``U자를 찾겠다''는 전제가
        \textbf{완전히 무너진다}.}
        \vskip0.4em
        \textbf{최적 지점을 찾는 실험은 반드시 여러 시드/반복으로
        평균내야 한다} --- 이 절이 실제로 강조하는 튜닝의
        전제 조건.
      \end{block}
    \end{column}
  \end{columns}
\end{frame}
